% !TEX root = ../main.tex

\chapter{Discussion}
\label{ch:discussion}

The evaluation in \cref{ch:results} measures the Boxel-based POD-TAMP
framework on three goals (holding, stack, find-and-tray-stack) against a
uniform-grid baseline that keeps everything else fixed, and against
TAMPURA's published numbers as an external reference. This chapter
interprets those results: where the adaptive partition pays off, where
finer leaves stop helping, what the TAMPURA comparison does and does not
show, and what the failure modes say about the bottleneck.

\section{Adaptive vs uniform partitioning}
\label{sec:disc-semantic-vs-uniform}

The headline finding is the cost gap between the adaptive partition and a
uniform free-space grid at the same cell scale. Both partitions use the
same auto-set cell size --- the largest object's footprint, $\approx$6--9\,cm
in our scenes --- so the comparison isolates the structural choice rather
than the resolution. On holding, the adaptive partition produces 25--43
Boxels per scene; the uniform grid produces $\sim$320 free-space cells
alone (\cref{fig:boxel-count-holding}). On stack the ratio widens to
$\sim$25 vs $\sim$1340. Larger partitions translate to larger initial-state
fact counts ($\sim$300 vs 11\,000--14\,000 on stack) and longer per-call
planning time, which is the mechanism behind \cref{fig:plantime-holding}.

What that gap buys is goal-dependent. On holding the uniform baseline
still solves about a third of scenes; it pays several times the planning
time at every $n_\text{occ}$ but the mechanism still works. On stack the
adaptive variants degrade gracefully from 97\,\% at $h=2$ to 13\,\% at
$h=4$, while uniform collapses from 97\,\% to 1\,\%. On
find-and-tray-stack, which composes occluded sensing with stacking,
uniform is effectively broken (1.3\,\% success). The reading is that the
adaptive partition does not change the planner; it changes how many
irrelevant cells the planner has to ground and search over. When the goal
needs only a handful of candidate placements --- a single hidden target ---
the overhead of a few hundred free-space cells is tolerable; when the
goal requires composing placements across multiple replans, the overhead
dominates.

\section{The 5\,cm leaf-floor variant}
\label{sec:disc-mbs0}

The \texttt{semantic+mbs0.05} variant was included to test whether
forcing a 5\,cm floor on the free-space leaf size changes outcomes
against the default \texttt{semantic} configuration, where the leaf
floor is the largest object's footprint. It does not: on stack the
curves overlap seed-for-seed; on holding the floor is $\sim$4\,pp ahead
of \texttt{semantic}; on find-and-tray-stack it is $\sim$7\,pp behind,
with a visible degradation as $n_\text{occ}$ grows
(\cref{fig:success-fats}). The mean replan count rises slightly on
holding and find-and-tray-stack relative to \texttt{semantic} ---
consistent with more sensing/replanning cycles without more
successes.

The mechanism is that the largest object's footprint ($\sim$6--9\,cm in
our object set) is already the binding constraint on placement; the 5\,cm
floor cannot shrink cells past this point and contributes only at the
cost of more cells in the partition. This is a characterisation of the
regime, not a defect of the method. In scenes whose placements are
constrained by sub-5\,cm geometry --- narrow pockets, dense small-occluder
shelves --- a finer leaf floor would in principle help. Constructing such
a scene and re-running the same comparison is the natural follow-up; it
is noted in \cref{sec:future_work}.

\section{Architectural comparison with TAMPURA}
\label{sec:disc-tampura}

\Cref{fig:tampura} reports our \texttt{semantic} median wall-clock of
14.0\,s on find-and-tray-stack against TAMPURA's reported
57\,$\pm$\,38\,s on its Partial Observability task (Table~II of
\cite{curtis2024partially}). The comparison merits a careful read
because of how the two systems differ.

The real difference is \emph{when} the geometry-sampling cost is paid.
TAMPURA pays it offline in a Learn-Model phase (Algorithm~2 of
\cite{curtis2024partially}): for each problem it constructs a sparse MDP
$\bar{\mathcal{B}}_\text{sparse}$ by running $K$ deterministic
FastDownward plans, then solves the sparse MDP at plan time with
LAO*. The expensive step --- sampling continuous parameters --- is done
once per problem and amortised across LAO*'s queries. Our system pays
the same kind of cost online, every PDDLStream call, with streams
sampling inside the planner loop. Both run a Fast Downward--based determinised planner inside the loop---ours
via PDDLStream; TAMPURA's released code via SymK, a Fast Downward--derived
top-$k$ planner (its Algorithm~2 names FastDownward)---so the salient
difference is \emph{when} geometry sampling is paid, not the search engine.
Neither system caches across episodes.

This framing changes what the 14.0\,s vs 57\,s number means. It is not a
``semantic-Boxels-beats-TAMPURA'' claim --- comparing wall-clocks of two
fundamentally different sampling schedules is not a benchmark of either.
It does show that the streams-online architecture is competitive with
the offline-MDP architecture on the closest analogous task, without front-loading
the sampling cost---and it reaches that parity without the learned sparse-MDP
abstraction the offline architecture must estimate, one of the representation
advantages set out in \cref{ch:related-work}: deterministic knowledge literals in
place of a probabilistic model, at the price of no risk-weighting and a reliance on
replanning (\cref{sec:limitations}).

Beyond \emph{when} sampling is paid, the comparison highlights \emph{what} each system can reason about. Because our planner holds occlusion as explicit symbolic state, line of sight is a planning condition: the \texttt{sense} action requires a clear view of its target region, and the objects that block a view are explicit facts in the planner's state. Those facts also cover candidate placements: a free cell lying on the camera's line of sight to a shadow is flagged as view-blocking, so the planner will not set an object down where it would block a region it still needs to observe. TAMPURA's \emph{Find Die} benchmark---its hidden-object task, the closest analogue to our setting---instead keeps occlusion sub-symbolic: a visibility voxel grid that only weights the success probability of the \texttt{look} action, and a placement sampler that draws object poses with no visibility check, so an object can be placed view-blind. Making space plannable in this way is the point of the representation; its present realization is partial---a clear view is a precondition, but \emph{restoring} a blocked view by first moving the obstruction is left to future work (\cref{sec:future_work}), and the current loop instead bounds itself with a give-up rule (\cref{sec:limitations}).

\section{Failure modes}
\label{sec:disc-failure}

\Cref{fig:failure-modes} stacks the exit reasons across all nine
$(\text{goal}, \text{variant})$ cells. Three patterns matter.

First, \texttt{planner\_failed} dominates everywhere. This means the
PDDLStream search exhausted its budget without finding a feasible plan
--- typically because every grasp or motion sample on a candidate Boxel
was rejected or the Boxel itself was unreachable. The bottleneck is
search, not perception or execution.

Second, \texttt{replan\_limit} is the second-largest slice on stack and
visible on holding \texttt{semantic+mbs0.05}. The sense--plan--act loop
is bounded by a maximum number of replans per episode; hitting this
bound indicates a shadow whose contents the system could not resolve
within the loop. This is the failure path of the bounded
three-strike give-up discussed in \cref{sec:limitations}: when a
sensing target returns \emph{still-blocked} three times the loop marks
the shadow empty and proceeds; if no candidate shadow remains the
episode fails with \texttt{replan\_limit}.

Third, \texttt{timeout} is small everywhere except find-and-tray-stack
adaptive variants, where it accounts for roughly 10\,\% of cells.
Find-and-tray-stack is a multi-step composition of search and stacking;
episodes that nearly succeed but exhaust the per-call budget on a
difficult replan show up here. \texttt{uniform} on find-and-tray-stack
is almost entirely \texttt{planner\_failed} (or \texttt{no\_summary},
which records the same fact at a different point in the execution
loop), consistent with the 1.3\,\% success rate.

Read together, the failure-mode breakdown points at search heuristics
and the give-up rule as the two paths most likely to lift success rates
without changing the representation.

\section{Threats to validity}
\label{sec:disc-validity}

The findings above are conditional on the scenes, the budget, and the
oracle perception. Five threats merit explicit mention.

\paragraph{Scene family.} All measurements are on tabletop scenes with a
fixed overhead camera, a Franka Panda, and 6--9\,cm cubes. The
adaptive-vs-uniform gap is most pronounced in this regime because
free-space cells dominate the uniform partition; in scenes with denser
objects or smaller workspaces the ratio would narrow.

\paragraph{Oracle perception.} Object detection reads ground-truth poses;
a learned detector consuming the rendered RGB-D would surface noise,
misclassification, and partial-detection failure modes the current
numbers do not reflect.

\paragraph{TAMPURA via published numbers.} We compare against TAMPURA's
Table~II rather than re-running their code. Algorithmic, implementation,
and parameter differences are bundled into a single delta. Both planners
are single-threaded, and TAMPURA's CPU has a marginally higher base clock
(2.5 vs 2.0\,GHz), so the hardware difference, if anything, slightly
favours TAMPURA rather than us. Re-running both systems on a shared scene
file would strengthen the comparison.

\paragraph{Resolution regime.} The free-space leaf size behaves as a flat
knob across the full $\sim$4$\times$ range tested. Below \emph{auto\_cell},
forcing a 5\,cm floor changes nothing (the \texttt{mbs0.05} null finding);
above it, coarsening the free-space leaf to 1.5$\times$ and 2$\times$
\emph{auto\_cell} leaves the success rate essentially unchanged --- identical
seed-for-seed on stack, within seed noise on holding, and only a mild
$\sim$4\,pp drop on find-and-tray-stack --- while shrinking the free-space
partition by 30--40\,\%. The cause is structural rather than numerical: the
partition's merge step collapses adjacent same-label free-space cells into a
single Boxel irrespective of the underlying leaf size, so leaf size alters
the partition only where sub-cell geometry forces a split. \emph{auto\_cell}
(the largest visible footprint plus a 1\,cm margin) is therefore a robust
default rather than a tuned constant, and semantic's compactness advantage
over \texttt{uniform} is a structural property of the adaptive partition, not
an artefact of resolution tuning. Neither result generalises to scenes where
placements are constrained by sub-5\,cm geometry; constructing such a scene
is the natural follow-up (\cref{sec:future_work}).

\paragraph{Bounded give-up.} The three-strike give-up on still-blocked
shadows is unsound for shadows that are physically unreachable. The run
report distinguishes shadows observed empty from shadows given up this
way, and the evaluation counts the give-up case as a failure;
\cref{sec:limitations} expands on the consequence.

The Limitations and Accepted Simplifications below catalogue the design
decisions that fix the scope of these results.

\section{Limitations and Accepted Simplifications}
\label{sec:limitations}

The system evaluated in this thesis embodies a set of deliberate
simplifications. Each was adopted to keep the implementation tractable and
to isolate the contribution --- semantic Boxel discretization and its
integration with partially observable deterministic TAMP --- from concerns
that are orthogonal to it. This subsection consolidates these
simplifications so that the empirical claims of \cref{ch:results} are
read with the correct scope.

\paragraph{Perception.} The scene contains a single fixed overhead camera
that observes the whole table. Its rendered RGB-D images are not processed,
however: object detection is performed by an oracle that queries
ground-truth object poses directly from the simulator and uses ray-casts to
determine which objects are visible from the camera's viewpoint. There is
therefore no learned detector and no sensor-noise model. This isolates the
planning contribution from perception error --- a real perception module
consuming the camera images could replace the oracle without altering the
planning architecture. Because the camera is fixed, the robot never moves
to a sensing configuration; this is sufficient for the tabletop scenes
studied here, where every shadow is visible from one viewpoint, but a
robot-mounted sensor would be required for occlusions not visible from any
fixed vantage point.

\paragraph{Manipulation and execution.} Grasping is constraint-based: a
held object is rigidly welded to the gripper, so objects cannot slip or be
dropped in transit. Only top-down grasps are sampled, at a single fixed
height offset. The final approach phase of a pick uses a local
inverse-kinematics solver that bypasses the planner's collision checks, on
the standard TAMP assumption that the short final motion is obstacle-free
if the pre-grasp pose is valid. After a place or stack, execution adds a
small hardcoded vertical lift before returning; this lift is invisible to
the planner --- it produces no PDDL state change and moves only along the
gravity axis --- and exists solely to give the sampling-based motion
planner a non-pathological start configuration for the next move. Collision
events detected during execution are logged but do not trigger an automatic
replan.

\paragraph{Spatial representation.} Free-space cells are merged only across
exactly aligned shared faces; a full semantic merge --- matching
view-blocking, observability, and ordering --- was not implemented, so the
partition contains more Boxels than strictly necessary. When an occluder is
relocated, its previously computed shadow geometry and the shadow-blocker
map are not recomputed at its new position. This rests on the assumption
that the world is static apart from the objects the robot itself moves:
nothing else shifts, so the only outdated geometry belongs to the relocated
object, whose new pose the planner tracks explicitly and replans from. The symbolic interface also uses
string identifiers as proxies for geometric volumes --- shadows and moved
occluders are passed to the planner as named entities rather than as
bounding boxes or occupancy grids --- bridging continuous geometry to
discrete planning at the cost of a representation that is not purely
geometric.

\paragraph{Planning.} The primary goal studied is \texttt{holding};
\texttt{stack} and \texttt{find-and-tray-stack} are shipped extensions that
broaden goal diversity. Enabling stacking roughly doubles per-call planning
time, traced to the conditional-effects requirement on the pick action;
this regression is accepted rather than optimized away, a mitigation having
been considered and dropped as out of scope. A second, smaller planner-search
bias concerns the cost model: the domain gives the \texttt{stack} action a cost
of 2 while every other action costs 1. This is a hand-tuned bias rather than a
claim about stacking---because the planner minimises total plan cost and cannot
weigh one action's outcome likelihood against another's, the higher cost is what
keeps it from substituting a \texttt{stack} as a cheap rescue when a \texttt{place}
onto a free Boxel fails its motion check, forcing it to exhaust placement options
first. It is a deterministic stand-in for the probabilistic action selection the
POD layer does not perform; \texttt{stack} remains available whenever the goal
requires it. More fundamentally, when a
shadow returns \emph{still-blocked} three times in succession the execution
loop marks it as not containing the target and proceeds. The shadow is
never directly observed, so for a physically unreachable shadow this
violates the invariant that belief should reflect observation. It is
accepted to keep the sense--plan--act loop bounded; the run report
distinguishes shadows observed empty from shadows given up this way, and
the evaluation counts the give-up case as a failure. The principled
remedy --- re-grounding the planner's view-blocking atoms from current
physics --- is left to future work (\cref{sec:future_work}).
\Cref{fig:retry-giveup} catches this event in the action log.

\begin{figure}[H]
    \centering
    \includegraphics[width=0.85\textwidth]{sim/sense_fail_retry3.png}
    \caption{The bounded give-up in action. The log line ``retry 3/3'' marks the third consecutive still-blocked sensing attempt on a shadow; at this point the loop marks the shadow empty and proceeds, even though it was never directly observed. This is the unsound shortcut discussed above.}
    \label{fig:retry-giveup}
\end{figure}

\paragraph{Parameters.} Several geometric constants --- grasp height
offsets, approach and lift heights, the octree's minimum leaf size, the
camera pose --- are tuned for the specific table, robot, and object set
used here and would not transfer unchanged to a different setup.
Generalizing them, for example by deriving grasp parameters from object
geometry, is orthogonal to the core contribution and is noted as future
work.

\paragraph{Simplifications disclosed in the approach.} Two further
simplifications are stated where they are introduced rather than here. The
conceptual belief model uses two Know-If literals, $K(p)$ and $K(\neg p)$;
the implemented domain collapses these into a single Know-If fluent paired
with a value-carrying fluent, an equivalent encoding for the scenarios
studied (\cref{sec:kif_belief}). The \texttt{sense} action is modeled
optimistically --- it assumes the target is found --- with reactive
replanning when execution reveals otherwise, because PDDLStream and
FastDownward do not support contingent planning (\cref{ch:methods}).
